\ticket{12}{}

\question{Формула Тейлора с остаточным членом в форме Пеано и в форме Лагранжа для ФНП}

\begin{note}
Идея: вводим $\varphi(t) = F(x_* + t\,dx)$ и сводим задачу к одномерному Тейлору.
\end{note}

\begin{theorem}[Остаток Лагранжа]
Пусть все частные производные $F$ до порядка $m+1$ существуют в $B_\delta(x_*)$.
Тогда $\forall x \in B_\delta(x_*)$, $dx = x - x_*$:
\[
    F(x) = \sum_{k=0}^{m} \frac{d^k_{x_*} F(dx)}{k!}
    + \frac{d^{m+1}_{x^\theta} F(dx)}{(m+1)!},
    \quad \theta \in (0,1),
    \quad x^\theta = x_* + \theta\,dx.
\]
\end{theorem}

\begin{proof}
Введём $\varphi(t) = F(x_* + t\,dx)$, тогда $\varphi(0) = F(x_*)$, $\varphi(1) = F(x)$.

\textit{Ключевой факт} (индукция по $k$): $\varphi^{(k)}(t) = d^k_{x^t}F(dx)$.

База $k=1$: внутренняя функция $H(t) = x_* + t\,dx$ имеет производную
$DH(t) = dx$, поэтому по теореме о дифференцировании композиции:
\[
    \varphi'(t) = DF(x^t)\cdot dx
    = \sum_{i=1}^n \frac{\partial F}{\partial x_i}(x^t)\cdot dx_i
    = d_{x^t}F(dx).
\]

Шаг: дифференцируем $d^l_{x^t}F(dx) =
\sum_{i_1,\ldots,i_l} \frac{\partial^l F}{\partial x_{i_1}\cdots\partial x_{i_l}}(x^t)
\cdot dx_{i_1}\cdots dx_{i_l}$ по $t$.
Множители $dx_{i_j}$ константны, поэтому снова применяем теорему о композиции
к каждой частной производной и получаем $d^{l+1}_{x^t}F(dx)$.

Итого $\varphi^{(k)}(0) = d^k_{x_*}F(dx)$.
Применяем одномерный Тейлор к $\varphi$ на $[0,1]$ и подставляем:
\[
    F(x) = \sum_{k=0}^{m} \frac{d^k_{x_*} F(dx)}{k!}
    + \frac{d^{m+1}_{x^\theta} F(dx)}{(m+1)!}. \qedhere
\]
\end{proof}

\begin{theorem}[Остаток Пеано]
Пусть все частные производные $F$ до порядка $m$ существуют и непрерывны
в $B_\delta(x_*)$. Тогда $\forall x \in B_\delta(x_*)$:
\[
    F(x) = \sum_{k=0}^{m} \frac{d^k_{x_*} F(dx)}{k!}
    + o\!\left(\|x - x_*\|^m\right), \quad x \to x_*.
\]
\end{theorem}

\begin{proof}
Применяем Лагранжа с $m-1$ и добавляем умный ноль $\pm\frac{d^m_{x_*}F(dx)}{m!}$:
\[
    F(x) = \sum_{k=0}^{m} \frac{d^k_{x_*} F(dx)}{k!} + R, \qquad
    R = \frac{d^m_{x^\theta}F(dx) - d^m_{x_*}F(dx)}{m!}.
\]

Оценим $R$. Так как $|dx_i| \leq \|x-x_*\|$, то
$|dx_{i_1}\cdots dx_{i_m}| \leq \|x-x_*\|^m$.
Обозначим:
\[
    \varepsilon(x) := \max_{i_1,\ldots,i_m}
    \left|
        \frac{\partial^m F}{\partial x_{i_1}\cdots\partial x_{i_m}}(x^\theta)
        -\frac{\partial^m F}{\partial x_{i_1}\cdots\partial x_{i_m}}(x_*)
    \right| \xrightarrow{x \to x_*} 0
\]
по непрерывности частных производных порядка $m$.
В сумме $n^m$ слагаемых, поэтому:
\[
    |R| \leq \frac{n^m \cdot \varepsilon(x)}{m!} \cdot \|x-x_*\|^m
    = o\!\left(\|x-x_*\|^m\right). \qedhere
\]
\end{proof}
\question{Интегральный признак сходимости ЧР}

% Ответ...

